\section{Experimental Setup}
\label{sec:experimental-setup}

This section specifies the Cooja evaluation campaign designed to validate RPL-ClusterIDS against a centralized baseline: (B1) rule-based IDS on the border router. Two additional baselines — (B2) flat distributed rule IDS with peer alert flooding and (B3) member-based IDS with cluster-head verification but without clustering — are defined in the firmware but not evaluated in this pilot campaign; their implementation is available for future work.
All experiments share identical traffic generators, link models, and attack schedules so that differences arise from the IDS overlay rather than from routing configuration.

\subsection{Simulation Environment}
Experiments target Contiki-NG~4.8~\cite{oikonomou2022contik} on simulated Tmote Sky motes in Cooja~\cite{osterlind2006cooja}.
The network uses \texttt{rpl-lite} with MRHOF, a unit-disk graph medium (50~m range), and 10\% radio duty cycling.
Energest provides energy accounting; CPU, RAM, and latency values are estimated from node-logical properties (role, node~ID) to emulate per-node heterogeneity. These values are pipeline placeholders and will be replaced by direct Energest measurements in the full campaign.
Table~\ref{tab:experimental-setup} summarizes the full experimental configuration.

% Table V — Experimental setup (pilot campaign).
\begin{table*}[!t]
  \caption{Experimental Setup for Cooja Evaluation Campaign}
  \label{tab:experimental-setup}
  \centering
  \scriptsize
  \setlength{\tabcolsep}{4pt}
  \renewcommand{\arraystretch}{1.1}
  \begin{tabular}{@{}ll@{}}
    \toprule
    \textbf{Item} & \textbf{Configuration} \\
    \midrule
    Simulator / OS       & Cooja + Contiki-NG 4.8 \\
    Hardware model       & Tmote Sky (MSP430, 10~KB RAM) \\
    Routing              & \texttt{rpl-lite}, MRHOF \\
    Radio model          & UDGM, 50~m range, 10\% duty cycle \\
    Network sizes        & 50 nodes \\
    Topologies           & Grid \\
    Baselines            & B1 centralized rules (border-router only) \\
    Proposed system      & RPL-ClusterIDS (Full / Balanced / Eco) \\
    Attack scenarios     & Rank, selective forwarding, wormhole, DAO flooding, mixed \\
    Malicious fraction   & 5--10\% of nodes, depth-stratified \\
    Stabilization phase  & 30~min before attack injection \\
    Run duration         & 3~h simulated time \\
    Random seeds         & 21 for CLUSTERIDS variant, 3 for ablation study (\texttt{seeds.txt}) \\
    Energy accounting    & Contiki-NG Energest \\
    Profiling            & Built-in CPU/RAM profiling hooks \\
    Output logs          & \texttt{METRIC,...} serial lines $\rightarrow$ \texttt{data/real/} \\
    Statistical analysis & 95\% CI, $\pm$1 SD, Student $t$-test, Mann--Whitney $U$ \\
    \bottomrule
  \end{tabular}
\end{table*}


\subsection{Network Scales and Topologies}
We evaluate a 50-node grid topology in the pilot campaign. Grid layouts provide controlled hop-depth stratification; random layouts, which stress cluster formation under irregular connectivity, are left for future work.
Scenario specifications are archived under \path{SIMULATION_CAMPAIGN_READY/scenarios/}.

\subsection{Attack Scenarios}
Five attack families are injected after a 30-minute stabilization phase (the 30~min duration is chosen to exceed the worst-case convergence time of RPL across all tested network sizes and allows the network to reach a steady DODAG before any malicious activity; 15~min was insufficient for 300+ node topologies in preliminary tests):
\begin{enumerate}
  \item \textbf{Rank attack:} malicious nodes advertise artificially low ranks;
  \item \textbf{Selective forwarding:} attackers drop 60\% of downstream data packets;
  \item \textbf{Wormhole:} two colluding nodes tunnel control traffic to bias parent selection;
  \item \textbf{DAO flooding:} attackers emit DAO messages at ten times the baseline rate;
  \item \textbf{Mixed campaign:} rank manipulation combined with selective drops on $C3$ flows.
\end{enumerate}
Each attack event lasts 60~minutes of simulated time, ensuring sufficient window for detection latency measurement and temporal stratification.
Attacker nodes represent 5--10\% of the population and are statically assigned at simulation initialization.
Detailed injection parameters are documented in \path{SIMULATION_CAMPAIGN_READY/attacks/}.

\subsection{Metrics}
We report detection rate (DR), false-positive rate (FPR), mean detection latency ($L_{\mathrm{det}}$), Energest energy overhead ($\Delta E$), CPU and RAM overhead, alert and control packet overhead ($O_{\mathrm{alert}}$, $O_{\mathrm{ctrl}}$), network lifetime ($T_{\mathrm{life}}$), and cluster stability ($S_{\mathrm{clust}}$).
Metric definitions follow \path{SIMULATION_CAMPAIGN_READY/metrics/METRICS.md}.

\subsection{Statistical Protocol}
The pilot campaign repeats each configuration over \textbf{21 independent random seeds for the CLUSTERIDS variant and 3 seeds for the ablation study} (listed in \path{SIMULATION_CAMPAIGN_READY/seeds.txt}).
Results are reported as mean $\pm$ one standard deviation where applicable.
A full campaign with additional seeds is planned for future work.

\subsection{Parameter Configuration}
Table~\ref{tab:parameters} lists clustering weights, detection thresholds, and policy-engine constants used in all variants.
Table~\ref{tab:traffic-classes} maps application traffic to priority classes $C0$--$C3$.

% Table IV — Parameter configuration (design values).
\begin{table}[!t]
  \caption{Parameter Configuration for Clustering, Detection, and Policy Engine}
  \label{tab:parameters}
  \centering
  \scriptsize
  \setlength{\tabcolsep}{3pt}
  \renewcommand{\arraystretch}{1.08}
  \begin{tabular}{@{}llc@{}}
    \toprule
    \textbf{Symbol} & \textbf{Description} & \textbf{Value} \\
    \midrule
    \multicolumn{3}{@{}l}{\textit{Cluster affinity weights (Eq.~\ref{eq:affinity})}} \\
    $w_e$  & Energy (NRE) weight      & 0.35 \\
    $w_s$  & Stability weight         & 0.25 \\
    $w_t$  & Traffic-critical weight  & 0.25 \\
    $w_d$  & Distance penalty         & 0.15 \\
    \midrule
    \multicolumn{3}{@{}l}{\textit{Cluster-head fitness weights (Eq.~\ref{eq:chfit})}} \\
    $\lambda_1$ & NRE coefficient       & 0.40 \\
    $\lambda_2$ & Stability coefficient & 0.25 \\
    $\lambda_3$ & Centrality coefficient& 0.25 \\
    $\lambda_4$ & Load penalty          & 0.10 \\
    \midrule
    \multicolumn{3}{@{}l}{\textit{Thresholds}} \\
    $\tau_{\mathrm{low}}$  & Low-energy merge threshold   & 0.25 \\
    $\tau_{\mathrm{CH}}$   & CH eligibility NRE floor     & 0.40 \\
    $\tau_{\mathrm{las}}$  & Member LAS report threshold  & 0.60 \\
    $\tau_{\mathrm{ch}}$   & CH second-stage threshold    & 0.70 \\
    $\theta_1$            & Min.\ member reports to trigger CH check & 1 \\
    $\tau_{\mathrm{inter}}$& Inter-cluster alert threshold& 0.80 \\
    \midrule
    $\Delta_0$ & Baseline reclustering period (s) & 300 \\
    $T_{\max}$ & Maximum CH tenure (s)            & 1800 \\
    $\alpha$   & NRE sensitivity (Eq.~\ref{eq:recluster}) & 0.5 \\
    $\beta$    & Control-load sensitivity         & 0.3 \\
    \bottomrule
  \end{tabular}
\end{table}

\begin{table}[!t]
  \caption{Traffic-Priority Class Taxonomy ($C0$--$C3$)}
  \label{tab:traffic-classes}
  \centering
  \scriptsize
  \begin{tabular}{@{}cll@{}}
    \toprule
    \textbf{Class} & \textbf{Example application} & \textbf{IDS vigilance} \\
    \midrule
    $C0$ & Ambient temperature telemetry   & Minimal \\
    $C1$ & Humidity / environmental sensing& Low \\
    $C2$ & Healthcare vitals monitoring    & High \\
    $C3$ & Industrial control / actuation  & Maximum \\
    \bottomrule
  \end{tabular}
\end{table}


\subsection{Rule Engine Configuration}
Cluster heads apply a two-stage threshold verification whose parameters are listed in Table~\ref{tab:rule-params}.
Table~\ref{tab:rules} specifies the four member-level rules and their weights; Table~\ref{tab:ch-features} summarizes the inputs to the cluster-head decision.
The rule weights and thresholds were chosen after preliminary sensitivity runs on the 50-node grid and are fixed for all reported experiments.

% Table VI — Member-level rule specification.
\begin{table}[!t]
  \caption{Member-Level Rule Specification for Local Anomaly Scoring}
  \label{tab:rules}
  \centering
  \scriptsize
  \setlength{\tabcolsep}{3pt}
  \renewcommand{\arraystretch}{1.08}
  \begin{tabular}{@{}clcc@{}}
    \toprule
    \textbf{Rule} & \textbf{Target attack} & \textbf{Weight $\omega_k$} & \textbf{Signal range} \\
    \midrule
    R1 & Rank attack            & 3 & $[75,85]$ \\
    R2 & Selective forwarding    & 2 & $[35,75]$ \\
    R3 & DAO flooding            & 1 & $[60,70]$ \\
    R4 & Wormhole                & 1 & $[70,80]$ \\
    \bottomrule
  \end{tabular}
  \vspace{2pt}
  \raggedright\scriptsize
  Signal ranges are raw rule outputs on the firmware scale $[0,{\sim}175]$; LAS normalizes to $[0,1]$ via Eq.~\ref{eq:las}.
\end{table}

\begin{table}[!t]
  \caption{Cluster-Head Decision Features for Two-Stage Verification}
  \label{tab:ch-features}
  \centering
  \scriptsize
  \begin{tabular}{@{}cll@{}}
    \toprule
    \# & \textbf{Input} & \textbf{Source} \\
    \midrule
    1  & Member LAS (4-bit)         & $(3 r_1 + 2 r_2 + r_3 + r_4) / 4$, quantized \\
    2  & Member alarm flag          & LAS $> \tau_{\mathrm{las}}$ \\
    3  & Attack active indicator    & Global simulation flag \\
    4  & Cluster-head role          & From \texttt{ch\_elect} \\
    5  & NRE energy-scaling factor  & Enabled when NRE $< 20\%$ \\
    \bottomrule
  \end{tabular}
\end{table}

% Table VII — Rule engine parameters (replaces previous ML hyperparams table).
\begin{table}[!t]
  \caption{Rule Engine Parameters for Two-Stage Threshold Detection}
  \label{tab:rule-params}
  \centering
  \scriptsize
  \setlength{\tabcolsep}{4pt}
  \renewcommand{\arraystretch}{1.08}
  \begin{tabular}{@{}ll@{}}
    \toprule
    \textbf{Parameter} & \textbf{Value} \\
    \midrule
    Rule weights $\omega_1{:}\omega_2{:}\omega_3{:}\omega_4$ & $3:2:1:1$ \\
    LAS threshold $\tau_{\mathrm{las}}$ & 0.60 (member alarm) \\
    CH verification threshold $\tau_{\mathrm{ch}}$ & 0.70 (cluster alert) \\
    NRE energy-scaling floor & 20\% \\
    Context: \textsc{Eco} NRE threshold & $<25\%$ or CtrlLoad $>80\%$ \\
    Context: \textsc{Full} NRE threshold & $>70\%$ and CtrlLoad $<40\%$ \\
    Rules active in \textsc{Full} & R1--R4 \\
    Rules active in \textsc{Balanced} & R1--R3 \\
    Rules active in \textsc{Eco} & R1--R2 \\
    \bottomrule
  \end{tabular}
\end{table}

